%% ma: L10-B02-0031
%% lop: 10
%% chuong: 1
%% bai: 2
%% dang: 10.2.2
%% loai: TN4
%% muc: TH
%% dap_an: "C"
%% nguon: "Ảnh giáo viên gửi 10/10/2026 (ThS. Nguyễn Hoàng Việt, Chương 2 Tập hợp), Đề 2 (đề thứ hai, không ghi tiêu đề), Phần 1 Câu 8"
%% kien_thuc: "Bài 2 (tập hợp và các phép toán trên tập hợp)"
%% trang_thai: cho-duyet
%% ly_do: "Dạng toán mới đề xuất 10.2.2: xác định tập hợp bằng cách liệt kê phần tử (giải phương trình, bất phương trình trong N, Z, Q, R)"
%% ghi_chu: "Đề gốc ghi Cho mệnh đề A=...; đã sửa thành Cho tập hợp A=... cho đúng"
\begin{ex}
Cho tập hợp $A=\left\{x\in\mathbb Q\mid(x+1)\left(x-\dfrac12\right)(x^2-2)=0\right\}$. Viết lại tập $A$ bằng phương pháp liệt kê.
\choice{$A=\left\{\dfrac12\right\}$}{$A=\{-1\}$}{\True $A=\left\{-1;\dfrac12\right\}$}{$A=\left\{-1;\dfrac12;\sqrt2;-\sqrt2\right\}$}
\loigiai{Các nghiệm: $x=-1$, $x=\dfrac12$, $x=\pm\sqrt2$. Vì $x\in\mathbb Q$ nên loại $\pm\sqrt2$ (vô tỉ). Vậy $A=\left\{-1;\dfrac12\right\}$.}
\end{ex}
